\uuid{S7tt}
\titre{Dénombrement}
\auteur{bodin ; hanani ; mzari}
\organisation{exo7}
\licence{CC BY-NC-SA 4.0}
\datecreate{2020-01-20}
\niveau{L1}
\module{Algèbre}
\chapitre{Logique, ensemble, raisonnement}
\sousChapitre{Ensemble}
\difficulte{3}
\theme{Ensembles, applications, dénombrement}
\exo7qcmid{lille-1.ensembles.029}

\contenu{
\begin{question}
\qclassification{100.02,101.01,102.01,102.02}
\qtags{motcle=dénombrement}
Soit $E$ un ensemble à $n$ éléments et $A\subset E$ une partie à $p < n$ éléments. On note $\mathcal{H}(E)$ l'ensemble des parties de $E$ qui contiennent un et un seul élément de $A$. Quel est le cardinal de $\mathcal{H}(E)$ ?
\begin{answers}
    \good{$\mathrm{Card}(\mathcal{H}(E))=p2^{n-p}$}
    \bad{$\mathrm{Card}(\mathcal{H}(E))=p$}
    \bad{$\mathrm{Card}(\mathcal{H}(E))=p2^p$}
    \bad{$\mathrm{Card}(\mathcal{H}(E))=p2^n$}
\end{answers}
\begin{explanations}
Si $A=\{a_1,\dots ,a_p\}$, les éléments de $\mathcal{H}(E)$ sont de la forme $\{a_i\}\cup B$, où $a_i\in A$ et $B\subset E\setminus A$. Donc $\mathrm{Card}(\mathcal{H}(E))=\mathrm{Card}(A)\times \mathrm{Card}(\mathcal{P}(E\setminus A))=p2^{n-p}$.
\end{explanations}
\end{question}
}
